\chapter{Introducción}

En la actualidad, el desarrollo de software se apoya en metodologías ágiles. En este enfoque, las historias de usuario se utilizan como medio de comunicación entre el equipo de desarrollo y las personas involucradas en el proyecto. Sin embargo, la redacción manual de estas historias requiere tiempo y esfuerzo, y puede introducir ambigüedades, omisiones o inconsistencias. Esto afecta la calidad de los requisitos y, en consecuencia, la del producto final.

En este contexto, el uso de técnicas de inteligencia artificial (IA), y en particular de modelos de lenguaje de gran escala (LLMs, por sus siglas en inglés), presenta el potencial de automatizar la generación y la validación de historias de usuario a partir de documentos de requisitos en distintos formatos. Esto puede agilizar el análisis, mejorar la trazabilidad entre los requisitos y las historias generadas y reducir errores asociados al trabajo manual. %Además, estas técnicas pueden procesar grandes volúmenes de información y relacionar fragmentos que aparecen distribuidos en uno o más documentos.

En este trabajo se busca comprender el estado actual de las herramientas que abordan este problema y evaluar las historias que generan. Para ello, se define un documento de requisitos, un conjunto de historias de usuario esperadas y estrategias para evaluar las historias generadas.

El documento se organiza en capítulos. En el Capítulo~\ref{cap:antecedentes} se presenta la revisión de antecedentes. El Capítulo~\ref{cap:herramientas} incluye la revisión de herramientas para la generación de historias de usuario, comenzando por la definición de la metodología de búsqueda, luego la obtención de los resultados, el filtrado y la selección de las herramientas. En el Capítulo~\ref{cap:central} se describen los pasos y decisiones tomadas tras la revisión de antecedentes, la evaluación piloto de una de las herramientas seleccionadas, ClickUp, la redefinición de los elementos de la evaluación y la propuesta de evaluación de historias de usuario. En el Capítulo~\ref{cap:evaluacion} se describe la metodología de evaluación, incluyendo los objetivos, los aspectos de calidad considerados y el procedimiento utilizado. En el Capítulo~\ref{cap:resultados} se presentan los resultados obtenidos a partir de la evaluación de la herramienta seleccionada, ClickUp. Finalmente, en el Capítulo~\ref{cap:conclusiones} se exponen las conclusiones del trabajo y posibles líneas de trabajo futuro.
